\subsection{Systems of equations in complete rings}

To abbreviate, we shall say in what follows that a ring satisfies Hensel's conditions if it is commutative, linearly topologized, Hausdorff and complete; given an ideal m in such a ring, m (or the ordered pair (A, m)) will be said to satisfy Hensel's conditions if m is closed in A and its elements are topologically nilpotent. The ideal t of A consisting of all the topologically nilpotent elements satisfies Hensel's conditions (no. 3).

In particular, if A is a commutative ring and m an ideal of A and A is Hausdorff and complete with respect to the m-adic topology, the ordered pair (A, m) satisfies Hensel's conditions.

PROPOSITION 6. Let A be a commutative ring, B a ring satisfying Hensel's conditions and \(u: \mathbf{A} \to \mathbf{B}\) a homomorphism. For every family \(\mathbf{x} = (x_1, \ldots, x_n)\) of topologically nilpotent elements of B, there exists a unique homomorphism \(\tilde{u}\) from \(\mathbf{A}[[\mathbf{X}_1, \ldots, \mathbf{X}_n]]\) to B such that \(\tilde{u}(a) = u(a)\) for all \(a \in \mathbf{A}\) and \(\tilde{u}(\mathbf{X}_i) = x_i\) for \(1 \leqslant i \leqslant n\) . Moreover, if m denotes the ideal of series without constant term in \(\mathbf{A}[[\mathbf{X}_1, \ldots, \mathbf{X}_n]]\) , \(\tilde{u}\) is continuous for the m-adic topology.

Let \(\mathfrak{a}\) be the finitely generated ideal generated in \(\mathbf{B}\) by the \(x_{i}\) ( \(1 \leqslant i \leqslant n\) ); for every open ideal \(\mathfrak{S}\) of \(\mathbf{B}\) , the images of the \(x_{i}\) in \(\mathbf{B}/\mathfrak{S}\) are nilpotent, hence the ideal \((\mathfrak{a} + \mathfrak{S})/\mathfrak{S}\) is nilpotent in \(\mathbf{B}/\mathfrak{S}\) and there exists an integer \(k\) such that, for \(\sum_{i=1}^{n} p_{i} \geqslant k\) , \(x_{1}^{p_{1}} \ldots x_{n}^{p_{n}} \in \mathfrak{S}\) . As every element of \(m^{k}\) is a finite sum of formal power series of the form \(X_{1}^{p_{1}} \ldots X_{n}^{p_{n}} g(X_{1}, \ldots, X_{n})\) , where \(\sum_{i=1}^{n} p_{i} \geqslant k\) , it is seen that, if \(\tilde{u}\) solves the problem, then \(\tilde{u}(\mathfrak{m}^k) \subset \mathfrak{S}\) , which proves the continuity of \(\tilde{u}\) . There obviously exists a unique homomorphism

\[
v \colon \mathrm{A} [ \mathrm{X} _ {1}, \dots , \mathrm{X} _ {n} ] \rightarrow \mathrm{B}
\]

such that \(v(a) = u(a)\) for \(a \in \mathbf{A}\) and \(v(\mathbf{X}_i) = x_i\) for \(1 \leqslant i \leqslant n\) and the above argument shows that \(v\) is continuous with respect to the topology induced on \(\mathbf{A}[\mathbf{X}_1, \ldots, \mathbf{X}_n]\) by the m-adic topology. As \(\mathbf{A}[\mathbf{X}_1, \ldots, \mathbf{X}_n]\) is dense in \(\mathbf{A}[[\mathbf{X}_1, \ldots, \mathbf{X}_n]]\) with the m-adic topology and B is Hausdorff and complete, this completes the proof of the existence and uniqueness of \(\tilde{u}\) .

Note that this proposition gives us again as a special case (i) of Proposition 11 of § 2, no. 9.

If A itself is linearly topologized, the restriction of \(\tilde{u}\) to \(A\{X_{1},\ldots,X_{n}\}\) coincides with the homomorphism derived from u in Proposition 4 of no. 2. This follows immediately from the fact that \(A[X_{1},\ldots,X_{n}]\) is dense in \(A\{X_{1},\ldots,X_{n}\}\) if this ring is given the topology with fundamental system of neighbourhoods of 0 the ideals \(m^{k}\cap N_{3}\) (in the notation of no. 2, this topology is the least upper bound of the topology induced on \(A\{X_{1},\ldots,X_{n}\}\) by the m-adic topology on \(A[[X_{1},\ldots,X_{n}]]\) and the topology defined in no. 2).

If B = A and u is the identity mapping, we shall write \(f(x_{1}, \ldots, x_{n})\) or \(f(\mathbf{x})\) for the element \(\tilde{u}(f)\) for every formal power series \(f \in A[[X_{1}, \ldots, X_{n}]]\) ; for every system \(\mathbf{f} = (f_{1}, \ldots, f_{r})\) of formal power series of \(A[[X_{1}, \ldots, X_{n}]]\) , let \(\mathbf{f}(\mathbf{x})\) denote the element \((f_{1}(\mathbf{x}), \ldots, f_{r}(\mathbf{x}))\) of \(A^{r}\) , then it is said to be obtained by substituting the \(x_{i}\) for the \(X_{i}\) in f.~If \(n \leqslant m\) and F is a formal power series of \(A[[X_{1}, \ldots, X_{m}]]\) , it is possible to consider F as a formal power series in \(X_{n+1}, \ldots, X_{m}\) with coefficients in \(A[[X_{1}, \ldots, X_{n}]]\) ; let

\[
\mathbf {F} (x _ {1}, \dots , x _ {n}, \mathbf {X} _ {n + 1}, \dots , \mathbf {X} _ {m})
\]

denote the formal power series in \(\mathrm{A}[[\mathrm{X}_{n+1}, \ldots, \mathrm{X}_{m}]]\) obtained by substituting the x \(_{i}\) for the X \(_{i}\) in the coefficients of F, for 1 ≤ i ≤ n.

Let us take B to be the ring of formal power series \(\Lambda[[\mathbf{X}_1,\ldots,\mathbf{X}_r]]\) and let \(n\) be the ideal of series in B without constant term, so that (B, n) satisfies Hensel's conditions (§ 2, no. 6, Corollary to Proposition 6). Proposition 6 may be applied by taking the \(x_i \in B\) to be series without constant term; then, for every series \(f \in A[[\mathbf{X}_1,\ldots,\mathbf{X}_n]]\) , \(\tilde{u}(f)\) is just the formal power series \(f(x_1,\ldots,x_n)\) defined in Algebra, Chapter IV, § 5, no. 5. This is obvious if \(f\) is a polynomial and it follows from the proposition in the general case by remarking that \(f \mapsto f(x_1,\ldots,x_n)\) is continuous on \(A[[\mathbf{X}_1,\ldots,\mathbf{X}_n]]\) with respect to the m-adic topology.

COROLLARY. Let A be a ring satisfying Hensel's condition and \(\mathbf{x} = (x_{1},\ldots ,x_{n})\) a family of topologically nilpotent elements of A. Let \(\mathbf{g} = (g_1,\dots,g_q)\) be a system without constant term of series in \(\mathbf{A}[[\mathbf{X}_1,\ldots ,\mathbf{X}_n]]\) and \(\mathbf{f} = (f_1,\dots,f_p)\) a system of formal power series in \(\mathbf{A}[[\mathbf{X}_1,\dots,\mathbf{X}_q]]\) . Then \(\mathbf{g}(\mathbf{x}) = (g_1(\mathbf{x}),\dots,g_q(\mathbf{x}))\) is a family of topologically nilpotent elements of \(\mathbf{A}\) and

\[
(\mathbf {f} \circ \mathbf {g}) (\mathbf {x}) = \mathbf {f} (\mathbf {g} (\mathbf {x})).\tag{18}
\]

The fact that the \(g_{i}(\mathbf{x})\) are topologically nilpotent follows immediately from Proposition 6 and the fact that in A the ideal of topologically nilpotent elements is closed. Relation (18) is obvious when the \(f_{j}\) are polynomials; on the other hand, if m and \(m'\) are the ideals of series without constant term in

\[
\mathrm{A} \left[ \left[ \mathrm{X} _ {1}, \dots , \mathrm{X} _ {q} \right] \right] \quad \text { and } \quad \mathrm{A} \left[ \left[ \mathrm{X} _ {1}, \dots , \mathrm{X} _ {n} \right] \right]
\]

respectively, clearly the relation \(f \in \mathfrak{m}^k\) implies \(f(g_1, \ldots, g_q) \in \mathfrak{m}'^k\) . The two sides of (18) are therefore continuous functions of \(\mathbf{f}\) to \((\mathrm{A}[[\mathrm{X}_1, \ldots, \mathrm{X}_q]])^p\) if \(\mathrm{A}[[\mathrm{X}_1, \ldots, \mathrm{X}_q]]\) is given the m-adic topology, by virtue of the above remark and Proposition 6; whence relation (18).

In what follows, for a ring A and an ideal m of A we shall denote by \(m^{\times n}\) the

product set \(\prod_{i=1}^{n} m_i\) in \(A^n\) , where \(m_i = m\) for \(1 \leqslant i \leqslant n\) , to avoid ambiguity.

PROPOSITION 7. Let A be a ring and m an ideal of A such that the ordered pair (A, m) satisfies Hensel's conditions. Let \(\mathbf{f} = (f_1, \ldots, f_n)\) be a system without constant term of series in \(\mathrm{A}[[\mathrm{X}_1, \ldots, \mathrm{X}_n]]\) such that \(\mathrm{J}_{\mathbf{f}}(0)\) is invertible in A. Then, for all \(\mathbf{x} \in \mathfrak{m}^{\times n}\) , \(\mathbf{f}(\mathbf{x}) \in \mathfrak{m}^{\times n}\) and \(\mathbf{x} \mapsto \mathbf{f}(\mathbf{x})\) is a bijection of \(\mathfrak{m}^{\times n}\) onto itself, the inverse bijection being \(\mathbf{x} \mapsto \mathbf{g}(\mathbf{x})\) , where \(\mathbf{g}\) is given by relation (16) in no. 4.

The fact that \(\mathbf{f}(\mathbf{x})\in\mathfrak{m}^{\times n}\) is obvious when the \(f_{i}\) are polynomials and follows in the general case from Proposition 6 and the fact that m is closed in A. The other assertions of the proposition are then immediate consequences of (16), (17) and (18).

COROLLARY. Let \(\mathfrak{q}\) be a closed ideal of A contained in \(\mathfrak{m}\) . Then the relation \(\mathbf{x} \equiv \mathbf{x}'\) (mod. \(\mathfrak{q}^{\times n}\) ) is equivalent to \(\mathbf{f}(\mathbf{x}) \equiv \mathbf{f}(\mathbf{x}')\) (mod. \(\mathfrak{q}^{\times n}\) ).

For every formal power series \(f \in \mathbf{A}[[\mathbf{X}_1, \ldots, \mathbf{X}_n]]\) ,

\[
f \left(\mathrm{X} _ {1}, \dots , \mathrm{X} _ {n}\right) - f \left(\mathrm{Y} _ {1}, \dots , \mathrm{Y} _ {n}\right) = \sum_ {i = 1} ^ {n} \left(\mathrm{X} _ {i} - \mathrm{Y} _ {i}\right) h _ {i} \left(\mathrm{X} _ {1}, \dots , \mathrm{X} _ {n}, \mathrm{Y} _ {1}, \dots , \mathrm{Y} _ {n}\right)
\]

where the \(h_i\) belong to \(\mathbf{A}[[\mathbf{X}_1,\ldots,\mathbf{X}_n,\mathbf{Y}_1,\ldots,\mathbf{Y}_n]]\) (Algebra, Chapter IV, § 5, no. 8, Proposition 9); it follows immediately that the relation \(\mathbf{x} \equiv \mathbf{x}'\) (mod. \(\mathfrak{q}^{\times n}\) ) implies \(\mathbf{f}(\mathbf{x}) \equiv \mathbf{f}(\mathbf{x}')\) (mod. \(\mathfrak{q}^{\times n}\) ). The converse is obtained by replacing \(\mathbf{f}\) by its ``inverse'' \(\mathbf{g}\) .

THEOREM 2. Let \(\mathbf{A}\) be a ring and \(\mathfrak{m}\) an ideal of \(\mathbf{A}\) such that the ordered pair \((\mathbf{A},\mathfrak{m})\) satisfies Hensel's conditions. Let \(\mathbf{f} = (f_1,\ldots ,f_n)\) be a system of \(n\) elements of \(\mathbf{A}\{\mathbf{X}_1,\ldots ,\mathbf{X}_n\}\) and let \(\mathbf{a}\in \mathbf{A}^n\) ; let us write \(\mathrm{J}_{\mathbf{f}}(\mathbf{a}) = e\) . There exists a system \(\mathbf{g} = (g_1,\dots ,g_n)\) of restricted formal power series without constant term in \(\mathbf{A}\{\mathbf{X}_1,\dots ,\mathbf{X}_n\}\) such that

\begin{enumerate}
\def\labelenumi{(\roman{enumi})}
\item
  \(M_{\mathbf{g}}(0) = I_{n}\) (unit matrix).
\item
  For all \(\mathbf{x} \in \mathbf{A}^n\) ,

\[
\mathbf {f} (\mathbf {a} + e \mathbf {x}) = \mathbf {f} (\mathbf {a}) + M _ {\mathbf {f}} (\mathbf {a}). e \mathbf {g} (\mathbf {x}).\tag{19}
\]

\item
  Let \(\mathbf{h} = (h_1, \ldots, h_n)\) be the system of formal power series without constant term (not necessarily restricted) such that \(\mathbf{g} \circ \mathbf{h} = \mathbf{1}_n\) (Proposition 5). For all \(y \in \mathfrak{m}^{\times n}\) ,

\[
\mathbf {f} (\mathbf {a} + e \mathbf {h} (\mathbf {y})) = \mathbf {f} (\mathbf {a}) + M _ {\mathbf {f}} (\mathbf {a}). e \mathbf {y}.\tag{20}
\]

\end{enumerate}

For every formal power series \(f \in \mathbf{A}[[\mathbf{X}_1, \ldots, \mathbf{X}_n]]\) ,

\[
f (\mathbf {X} + \mathbf {Y}) = f (\mathbf {X}) + M _ {f} (\mathbf {X}). \mathbf {Y} + \sum_ {1 \leqslant i \leqslant j \leqslant n} \mathrm{G} _ {i j} (\mathbf {X}, \mathbf {Y}) \mathrm{Y} _ {i} \mathrm{Y} _ {j}\tag{21}
\]

where the \(G_{ij}\) are well determined formal power series in

\[
\mathrm{A} \left[ \left[ \mathrm{X} _ {1}, \dots , \mathrm{X} _ {n}, \mathrm{Y} _ {1}, \dots , \mathrm{Y} _ {n} \right] \right].
\]

If \(f\) is restricted, so are the elements of \(M_{f}\) and the \(G_{ij}\) , for these formal power series are polynomials if \(f\) is a polynomial and it follows from their uniqueness that for every open ideal \(\Im\) of \(A\) , denoting by \(p_{\Im} \colon A \to A / \Im\) the canonical mapping, the image of \(G_{ij}\) under \(\bar{p}_{\Im}\) is the coefficient of \(Y_{i}Y_{j}\) in \(\bar{p}_{\Im}(F)\) where \(F\) is the formal power series \(f(\mathbf{X} + \mathbf{Y})\) in \(A[[X_1, \ldots, X_n, Y_1, \ldots, Y_n]]\) ; whence our assertion.

This being so, writing formula (21) for each series \(f_{i}\) ( \(1 \leqslant i \leqslant n\) ), we obtain for all \(\mathbf{x} \in \mathbf{A}^n\) (no. 2, Proposition 4),

\[
\mathbf {f} (\mathbf {a} + e \mathbf {x}) = \mathbf {f} (\mathbf {a}) + M _ {\mathbf {f}} (\mathbf {a}). e \mathbf {x} + e ^ {2} \mathbf {r} (\mathbf {x})\tag{22}
\]

where \(\mathbf{r} = (r_1, \ldots, r_n)\) is a system of restricted formal power series each of which is of total order \(\geqslant 2\) . It follows from formulae (18) of Algebra, Chapter III, §6, no. 5 that there exists a square matrix \(M' \in \mathbf{M}_n(\mathbf{A})\) such that

\[
M _ {\mathbf {f}} (\mathbf {a}). M ^ {\prime} = e I _ {n},\tag{23}
\]

whence using this in (22)

\[
\mathbf {f} (\mathbf {a} + e \mathbf {x}) = \mathbf {f} (\mathbf {a}) + M _ {\mathbf {f}} (\mathbf {a}). e \mathbf {x} + M _ {\mathbf {f}} (\mathbf {a}) M ^ {\prime}. e \mathbf {r} (\mathbf {x}).\tag{24}
\]

Writing \(\mathbf{g} = \mathbf{1}_n + M'.\mathbf{r}\) , we see that \(\mathbf{g}\) satisfies conditions (i) and (ii); then it is sufficient to replace \(\mathbf{x}\) by \(\mathbf{h}(\mathbf{y})\) to obtain (iii).

COROLLARY 1. Let A be a ring and m an ideal of A such that the ordered pair (A, m) satisfies Hensel's conditions. Let \(f \in \mathrm{A}\{\mathbf{X}\}\) , \(a \in \mathrm{A}\) and write \(e = f'(a)\) . If \(f(a) \equiv 0 \pmod{e^2\mathfrak{m}}\) , then there exists \(b \in \mathrm{A}\) such that \(f(b) = 0\) and \(b \equiv a \pmod{e\mathfrak{m}}\) . If \(b'\) is another element of A such that \(f(b') = 0\) and \(b' \equiv a (\text{mod. em})\) , then \(e(b - b') = 0\) . In particular, b is unique if e is not a divisor of zero in A.

Let \(f(a) = e^{2}c\) where \(c \in \mathfrak{m}\) ; formula (20) for \(n = 1\) gives

\[
f (a + e h (y)) = e ^ {2} (c + y)
\]

and it is therefore sufficient to take \(y = -c\) , \(b = a + eh(-c)\) . Moreover if \(b = a + ex\) , \(b' = a + ex'\) , \(x \in \mathfrak{m}\) , \(x' \in \mathfrak{m}\) , \(f(b) = f(b') = 0\) , we deduce from (19) that \(e^2(g(x) - g(x')) = 0\) . As \(g(\mathbf{X}) - g(\mathbf{Y}) = (\mathbf{X} - \mathbf{Y})u(\mathbf{X}, \mathbf{Y})\) , where \(u\) is restricted and \(u(0,0) = 1\) , \(g(x) - g(x') = (x - x')v\) , where \(v \in \mathbf{A}\) is invertible, for, since \(\mathfrak{m}\) is closed, \(v - 1 = u(x,x') - 1 \in \mathfrak{m}\) and \(\mathfrak{m}\) is contained in the Jacobson radical of \(\mathbf{A}\) ; whence the relation \(e(b - b') = 0\) .

Remark. The corollary applies notably when \(e\) is invertible in A; we can then also deduce the existence of \(b\) from Hensel's Theorem, for the canonical image of \(f(\mathbf{X})\) in \((\mathbf{A} / \mathfrak{m})\{\mathbf{X}\}\) is of the form \((\mathbf{X} - \alpha)f_1(\mathbf{X}), \mathbf{X} - \alpha\) and \(f_1(\mathbf{X})\) being strongly relatively prime, for \(f_1(\alpha) = f'(\alpha)\) is the image of \(e\) (no. 1, Example).

\textbf{Examples}\par

\begin{enumerate}
\def\labelenumi{(\arabic{enumi})}
\item
  Let \(p\) be a prime number \(\neq 2\) and \(n\) an integer whose class mod. \(p\) is a square \(\neq 0\) in the prime field \(\mathbf{F}_p\) . If \(\mathbf{Z}_p\) is the ring of \(p\) -adic integers (§ 2, no. 12, Example 3), the application of Corollary 1 to the polynomial \(\mathbf{X}^2 - n\) shows that \(n\) is a square in \(\mathbf{Z}_p\) ; for example 7 is a square in \(\mathbf{Z}_3\) .
\item
  Let \(\mathbf{A} = \mathbf{K}[[\mathbf{Y}]]\) be the ring of formal power series in one indeterminate with coefficients in a commutative field \(\mathbf{K}\) ; with the (Y)-adic topology, the ring \(\mathbf{A}\) is Hausdorff and complete (§ 2, no. 6, Corollary to Proposition 6) and the mapping \(f(\mathbf{Y}) \mapsto f(0)\) defines by passing to the quotient ring an isomorphism of \(\mathbf{A}/(\mathbf{Y})\) onto the field \(\mathbf{K}\) . By Corollary 1, if \(\mathbf{F}(\mathbf{Y}, \mathbf{X})\) is a polynomial in \(\mathbf{X}\) with coefficients in \(\mathbf{A}\) and \(a\) is a simple root of \(\mathbf{F}(0, \mathbf{X})\) in \(\mathbf{K}\) , there exists a unique formal power series \(f(\mathbf{Y})\) such that \(f(0) = a\) and \(\mathbf{F}(\mathbf{Y}, f(\mathbf{Y})) = 0\) .
\end{enumerate}

COROLLARY 2. Let A be a ring and m an ideal of A such that the ordered pair (A, m) satisfies Hensel's conditions. Let r, n be integers such that \(0 \leqslant r < n\) and \(\mathbf{f} = (f_{r+1}, \ldots, f_n)\) is a system of n - r elements of \(A\{X_1, \ldots, X_n\}\) ; let \(\mathrm{J}_{\mathbf{f}}^{(n-r)}(\mathbf{X})\) denote the minor of \(M_{\mathbf{f}}(\mathbf{X})\) consisting of the columns of index j such that \(r + 1 \leqslant j \leqslant n\) . Let \(a \in A^n\) be such that \(\mathrm{J}_{\mathbf{f}}^{(n-r)}(\mathbf{a})\) is invertible in A and \(\mathbf{f}(\mathbf{a}) \equiv 0 \pmod{\mathfrak{m}^{\times(n-r)}}\) . Then there exists a unique \(\mathbf{x} = (x_1, \ldots, x_n) \in \mathrm{A}^n\) such that \(x_k = a_k\) for \(1 \leqslant k \leqslant r\) , \(x \equiv a \pmod{m^{\times n}}\) and \(\mathbf{f}(\mathbf{x}) = 0\) .

Substituting \(a_k\) for \(\mathbf{X}_k\) for \(1 \leqslant k \leqslant r\) in the \(f_i\) (no. 2, Remark 3), we see immediately that we may restrict our attention to the case where \(r = 0\) to prove the corollary. Theorem 1 and Proposition 7 show then that \(\mathbf{f}\) defines a bijection of \(\mathbf{a} + \mathfrak{m}^{\times n}\) onto \(\mathbf{f}(\mathbf{a}) + \mathfrak{m}^{\times n} = \mathfrak{m}^{\times n}\) ; the corollary follows from the fact that \(0 \in \mathfrak{m}^{\times n}\) .

COROLLARY 3. With the notation of Corollary 2, let \(\mathbf{a} \in \mathbf{A}^n\) ; let us write \(e = \mathrm{J}_{\mathbf{f}}^{(n-r)}(\mathbf{a})\) (not necessarily invertible in A) and suppose that \(\mathbf{f}(\mathbf{a}) \equiv 0 \pmod{e^2\mathfrak{m}^{\times(n-r)}}\) . Then there exist \(n - r\) formal power series without constant term \(\phi_i (r + 1 \leqslant i \leqslant n)\) in \(\mathbf{A}[[\mathbf{X}_1, \ldots, \mathbf{X}_r]]\) such that, for all \(\mathbf{t} = (t_1, \ldots, t_r) \in \mathfrak{m}^{\times r}\) ,

\[
f _ {i} \left(a _ {1} + e ^ {2} t _ {1}, \dots , a _ {r} + e ^ {2} t _ {r}, a _ {r + 1} + e \phi_ {r + 1} (\mathbf {t}), \dots , a _ {n} + e \phi_ {n} (\mathbf {t})\right) = 0
\]

\[
\text{for } r + 1 \leqslant i \leqslant n.
\]

For \(1 \leqslant i \leqslant r\) , let \(f_{i}(\mathbf{X}) = \mathbf{X}_{i} - a_{i}\) and let \(\mathbf{u} = (f_{1}, \ldots, f_{n})\) ; then \(\mathrm{J}_{\mathbf{u}}(\mathbf{a}) = e\) and Theorem 2 may be applied to the system \(\mathbf{u}\) . With the notation of Theorem 2 it follows from the above definitions that \(g_{i}(\mathbf{X}) = \mathbf{X}_{i}\) for \(1 \leqslant i \leqslant r\) , whence \(h_{i}(\mathbf{X}) = \mathbf{X}_{i}\) for \(1 \leqslant i \leqslant r\) ; moreover, if \(M' \in \mathbf{M}_{n}(\mathbf{A})\) is such that \(M_{\mathbf{u}}(\mathbf{a}) \cdot M' = eI_{n}\) , \(M'\) is of the form

\[
\left( \begin{array}{c c} e I _ {r} & 0 \\ * & * \end{array} \right).
\]

Replacing \(\mathbf{y}\) by \(M^{\prime}.\mathbf{z}\) (where \(\mathbf{z} = (z_{1},\ldots ,z_{n})\in \mathfrak{m}^{\times n}\) ) in formula (20), we obtain

\[
\begin{array}{r l} f _ {i} (a _ {1} + e ^ {2} z _ {1}, \dots , a _ {r} + e ^ {2} z _ {r}, a _ {r + 1} + e h _ {r + 1} (M ^ {\prime}. \mathbf {z}), \dots , a _ {n} + e h _ {n} (M ^ {\prime}. \mathbf {z})) \\ & = f _ {i} (a) + e ^ {2} z _ {i} \quad \text { for } \quad 1 \leqslant i \leqslant n. \end{array}\tag{26}
\]

By hypothesis, \(f_{j}(a) = e^{2}b_{j}\) where \(b_{j} \in \mathfrak{m}\) for \(r + 1 \leqslant j \leqslant n\) . Let us write \(\psi_{j}(\mathbf{X}_{1},\ldots ,\mathbf{X}_{n}) = h_{j}(M^{\prime}.\mathbf{X})\) and

\[
\phi_ {j} \left(\mathrm{X} _ {1}, \dots , \mathrm{X} _ {r}\right) = \psi_ {j} \left(\mathrm{X} _ {1}, \dots , \mathrm{X} _ {r}, - b _ {r + 1}, \dots , - b _ {n}\right)
\]

for \(r + 1 \leqslant j \leqslant n\) . For \(r + 1 \leqslant i \leqslant n\) , substituting \(t_j\) for \(z_j\) for \(1 \leqslant j \leqslant r\) and \(b_j\) for \(z_j\) for \(r + 1 \leqslant j \leqslant n\) in (26), we obtain relations (25) for all \(t \in m^{\times r}\) .

